\subsection{From Saturday evening to Sunday evening}

The first time the Church tells us to get up it is still Saturday
evening. At First Vespers the short reading is St Paul's summons to
wake from sleep. The hymn that follows addresses Jesus as Creator of
the stars. It remembers how far he came to heal the world. From his
Virgin Mother he went to the Cross, and the singers ask the one who
will judge them to defend them now. In Cardinal Newman's English the
Church sings to the Lord who came “to be the medicine / Of a sick
world, in love”
(\emph{Creator alme siderum}, Saturday before the First Sunday).
It is possible to need this medicine and to fear the very one who can
give it. From its first evening Advent teaches us to put both things
into the same prayer.

An antiphon is a short chant sung with a psalm or canticle. The
antiphon for Mary's Magnificat this evening announces that the name of
the Lord comes from afar and that his brightness fills the whole
earth. The woman whose words the choir then sings has felt the reach
of that brightness in her own life. She begins with what the Almighty
has done for her. By the end of her song she is speaking for Israel,
remembering the promise to Abraham and his descendants (Luke
1:46--55).

Matins is the Hour with the long readings. On the First Sunday the
prophet Isaiah opens by summoning heaven and earth to listen to a
father speak. He has fed and brought up children who have turned
against him. An ox knows its owner and a donkey knows the manger
where its master feeds it. The people God has fed have failed to
recognize him (Isa.~1:1--3). There is a question to take from the
first night of Advent to the crib on Christmas night. Do I recognize
the one who feeds me? When someone reads this passage from the
Breviary and later stands beside a manger scene, even the place
where the child lies can call Isaiah's question back to mind.

Someone using an older English Breviary to follow the readings may
find nine lessons in Bute's book. In the Office of an Advent Sunday
under the 1960 rubrics there are three. Two come from Scripture.
Then the opening of the Gospel and part of a homily are read.
The first lesson of Bute's old third nocturn, the third group of
readings in the night Office, supplies that homily. On the First
Sunday St Gregory opens Luke's account of the last day by asking
what the Saviour wishes to find when he comes. The judge himself
wants to find his people ready. The Hour ends its readings with a
response in place of the \emph{Te Deum}, the solemn hymn of praise
(\emph{General Rubrics of the Breviary}, nos.~162, 166, 220).

At Mass the singers ask God to show his mercy and grant his salvation.
Those are the words of Psalm 84 sung with the Alleluia before the
First Sunday's Gospel. The Gospel speaks of the coming of the Son of
Man. At Communion the same psalm is heard again, this time with a
promise that the Lord will give what is good and our earth will
yield its fruit (Ps.~84:8, 13). A person who heard the petition on
the way to the Gospel can hear the promise on the way to the altar.
God's answer to those who await the Saviour includes the gift of his
Son to them here.

The day has still more of Paul's letter to give. The short reading
at Sext, the Hour of midday, says that night has passed and the day
is near. At None the reading names some of the things to give up,
including quarrels and envy, and bids us put on the Lord Jesus
Christ (Rom.~13:11--14). The same Epistle holds together the eve,
the Sunday Mass and the afternoon. Someone who has found his temper
shortened by midday has another chance, a few hours later, to hear
what it means to dress for the day that is coming.

At Second Vespers the antiphon for the Magnificat changes. The
first evening had announced light reaching the ends of the earth.
On Sunday evening the choir speaks to one woman. It takes up
Gabriel's words telling Mary to put away fear, for she has found
favour and will bear a son (Luke 1:30--31). Then her own canticle is
sung again. Within one Sunday the Church has heard that nations will
be in distress when the Son of Man appears, and has sung to the
woman who first received him. The voice Mary had heard at Nazareth
gave her reason to trust the Lord who comes.

\subsection{What the weekdays carry forward}

On a weekday without a Mass of its own the Missal directs the priest
celebrating the Mass of the feria to return to the preceding Sunday's
Mass. A \emph{feria} here is a weekday of the season. On each of the
four Sundays of Advent an instruction stands between the Gradual,
the psalm chant after the Epistle, and the Alleluia. At a weekday
celebration that uses this Mass the Gradual is retained and the
Alleluia and its verse are left out. On a weekday that takes the
First Sunday's Mass, the last words sung before the Gospel thus ask
the Lord to show us his ways and teach us his paths (Ps.~24:3--4).
A listener can bring a decision to those words and hear the Gospel
as someone who has just asked to be taught.

Some weekdays will give the reader a new Mass altogether. The Ember
Days tell of the angel coming to Mary, Mary going to Elizabeth, and
the voice that prepares a way in the wilderness. The Office can also
move forward when the Mass returns to a Sunday. On the evenings
before Christmas the Magnificat of the Advent Office has an
antiphon for the date. On December 17 the singers call on Wisdom to
come and teach them. By December 23 they are addressing Emmanuel,
God with us. Mary's song stays with them while the words they bring
to it change from evening to evening. It gives a household a way to
keep Advent together even when its members can meet for prayer only
at the end of the day.
